\documentclass[11pt]{article}

% arXiv-friendly packages
\usepackage[utf8]{inputenc}
\usepackage[T1]{fontenc}
\usepackage{hyperref}
\usepackage{url}
\usepackage{booktabs}
\usepackage{amsmath}
\usepackage{graphicx}
\usepackage[margin=1in]{geometry}
\usepackage{natbib}
\usepackage{xcolor}
\usepackage{listings}

\lstset{
  basicstyle=\small\ttfamily,
  frame=single,
  breaklines=true,
  columns=flexible
}

\title{IMI: Integrated Memory Intelligence for AI Agents\\
\large A Cognitive Memory Architecture with Benchmarks}

\author{
  Renato Aparecido Gomes \\
  \texttt{renatoapgomes@gmail.com} \\
  \url{https://github.com/RAG7782/imi}
}

\date{March 2026}

\begin{document}
\maketitle

\begin{abstract}
We present IMI (Integrated Memory Intelligence), a memory system for AI agents
that goes beyond vector-based retrieval by integrating cognitive features:
affective modulation, temporal decay, surprise-based encoding, affordance
extraction, multi-resolution zoom, and graph-augmented multi-hop retrieval.
Through systematic ablation (100 memories, 15 queries), we show that while
cognitive features trade off pure retrieval quality (R@5 drops from 0.341 to
0.204 at relevance\_weight=0.3), they significantly improve agent-relevant
metrics: domain precision increases +6.7\% for temporal queries, and
graph-augmented expansion achieves 100\% multi-hop recall (vs 75\% cosine-only).
We also introduce AMBench, the first standardized benchmark for agent memory
systems, evaluating the full encode-retrieve-consolidate-act lifecycle over
90 simulated days with 300 incidents. IMI is open-source, requires no external
infrastructure (SQLite-only), and operates with zero LLM calls at query time.
\end{abstract}

\section{Introduction}

Current AI memory systems face a fundamental tension: they optimize for
\textbf{retrieval accuracy} (finding the most similar document) when agents
actually need \textbf{decision relevance} (finding the most useful memory for
the current situation).

RAG systems use cosine similarity over dense embeddings. This works for
knowledge Q\&A but fails for agent memory, which requires:

\begin{itemize}
  \item \textbf{Temporal coherence}: Recent incidents are often more relevant than old ones
  \item \textbf{Action orientation}: ``What can I DO?'' not just ``What do I KNOW?''
  \item \textbf{Multi-hop reasoning}: Causal chains span multiple memories
  \item \textbf{Emotional salience}: Critical incidents should be harder to forget
  \item \textbf{Adaptive resolution}: Sometimes you need the gist, sometimes the details
\end{itemize}

IMI addresses these gaps through a cognitive architecture inspired by human
hippocampal memory, ecological psychology \citep{gibson1979}, and predictive
coding theory \citep{rao1999}.

\subsection{Contributions}

\begin{enumerate}
  \item \textbf{IMI Architecture}: A cognitive memory system with 7 composable features
    (affect, surprise, affordances, temporal decay, mass, zoom, graph edges)
  \item \textbf{Ablation Study}: Systematic measurement of each feature's contribution
    to retrieval quality and agent relevance
  \item \textbf{AMBench}: First standardized benchmark for AI agent memory (5 metrics,
    3 baselines, 300 incidents over 90 simulated days)
  \item \textbf{Graph-Augmented Retrieval}: Lightweight edge layer that achieves 100\%
    multi-hop recall with zero LLM calls
  \item \textbf{Empirical finding}: The optimal relevance weight for agent memory is
    0.10--0.15, not the 0.3 commonly assumed in prior work
\end{enumerate}

\section{Related Work}

\begin{table}[h]
\centering
\small
\begin{tabular}{lccccc}
\toprule
System & Retrieval & Temporal & Affordances & Multi-hop & Benchmark \\
\midrule
RAG (vanilla) & Cosine & No & No & No & No \\
HippoRAG (2024) & KG + PPR & No & No & Yes & No \\
RAPTOR (2024) & Tree & No & No & Partial & No \\
GraphRAG (2024) & Entity KG & No & No & Yes & No \\
MemGPT (2023) & Tiered & Yes & No & No & No \\
\textbf{IMI (ours)} & \textbf{Cos+Graph} & \textbf{Yes} & \textbf{Yes} & \textbf{Yes} & \textbf{Yes} \\
\bottomrule
\end{tabular}
\caption{Comparison of memory systems for AI agents.}
\label{tab:related}
\end{table}

\textbf{HippoRAG} \citep{gutierrez2024} is the closest to IMI. Both are
hippocampus-inspired, but HippoRAG focuses on knowledge retrieval via
Personalized PageRank over a knowledge graph built with LLM NER. IMI focuses
on agent memory with temporal decay, affordances, and zero LLM calls at
query time.

\textbf{MemGPT} \citep{packer2023} implements tiered memory (working/long-term)
with LLM-controlled paging. IMI's zoom levels serve a similar purpose but
are query-time, not storage-time.

\textbf{RAPTOR} \citep{sarthi2024} uses recursive tree summarization for
multi-resolution retrieval. IMI's zoom levels (orbital/medium/detailed) are
analogous but pre-computed at encode time.

\section{Architecture}

\subsection{Memory Node}

Each memory is a \texttt{MemoryNode} with:
\begin{itemize}
  \item \textbf{seed}: original experience text
  \item \textbf{embedding}: 384-dimensional vector (all-MiniLM-L6-v2)
  \item \textbf{summaries}: orbital (gist), medium, detailed (3 zoom levels)
  \item \textbf{affect}: AffectiveTag (salience, valence, arousal) $\rightarrow$ modulates decay
  \item \textbf{affordances}: list of action potentials (``what can I do with this?'')
  \item \textbf{mass}: gravitational pull (derived from affect encoding strength)
  \item \textbf{tags}: domain labels for co-occurrence linking
\end{itemize}

\subsection{Relevance Scoring}

The relevance score combines temporal, frequency, mass, and surprise signals:

\begin{equation}
\text{relevance} = \text{recency} \times (1 + \text{frequency}) \times \text{mass} \times \text{surprise\_boost}
\end{equation}

Where:
\begin{align}
\text{recency} &= \frac{1}{1 + \text{days\_since} \times (1 - 0.5 \times \text{fade\_resistance})} \\
\text{frequency} &= \log(1 + \text{access\_count}) \\
\text{mass} &= \text{affect.encoding\_strength}
\end{align}

Final search score:
\begin{equation}
\text{score} = (1 - rw) \times \cos(\mathbf{q}, \mathbf{m}) + rw \times \hat{r}
\end{equation}

Optimal $rw = 0.10$ (validated via ablation and temporal experiments).

\subsection{Graph Layer}

Lightweight edge layer over the vector store:
\begin{itemize}
  \item \textbf{Edge types}: causal, co\_occurrence, similar
  \item \textbf{Auto-linking}: $\cos > 0.75$ triggers similar edge (no LLM calls)
  \item \textbf{Retrieval}: cosine seeds $\rightarrow$ graph expansion (spreading activation) $\rightarrow$ re-rank
\end{itemize}

Score with graph:
\begin{equation}
\text{score} = (1-rw-gw) \times \cos + rw \times \hat{r} + gw \times \text{activation}
\end{equation}

\section{Experiments}

\subsection{Dataset: 100 SRE Postmortems}

100 synthetic incident reports across 5 domains (auth, database,
infrastructure, monitoring, network), 20 per domain, with known causal chains
and ground truth relevance judgments.

\subsection{Ablation Study (WS-A)}

\begin{table}[h]
\centering
\begin{tabular}{lrrl}
\toprule
Feature Removed & $\Delta$R@5 & $\Delta$MRR & Effect \\
\midrule
All features ($rw$=0$\rightarrow$0.3) & $-$0.137 & $-$0.126 & Hurts retrieval \\
Recency & $-$0.117 & $-$0.117 & Hurts most \\
Affect + Mass & $-$0.107 & $-$0.103 & Hurts \\
Mass only & $-$0.108 & $-$0.099 & Hurts \\
Frequency & $-$0.061 & $-$0.069 & Hurts \\
Surprise & $-$0.003 & $-$0.004 & Negligible \\
\bottomrule
\end{tabular}
\caption{Ablation study: effect of removing each feature on retrieval metrics.
Features that model human memory trade retrieval accuracy for agent relevance.}
\label{tab:ablation}
\end{table}

\subsection{Temporal Decay Test (WS-B)}

With 90 days of simulated access patterns (power-law distribution):

\begin{table}[h]
\centering
\begin{tabular}{lrrr}
\toprule
Scenario & $rw$=0.0 & $rw$=0.15 & $\Delta$ \\
\midrule
Overall domain P@5 & 0.689 & 0.756 & +0.067 \\
Recent/recurring & 0.800 & 0.900 & +0.100 \\
Old/forgotten & 0.600 & 0.600 & 0.000 \\
\bottomrule
\end{tabular}
\caption{Relevance weighting improves domain precision +6.7\% without
degrading access to old memories.}
\label{tab:temporal}
\end{table}

\subsection{Graph-Augmented Multi-hop (WS-G)}

\begin{table}[h]
\centering
\begin{tabular}{lrrl}
\toprule
System & Multi-hop R@10 & Hits & Std R@5 \\
\midrule
Cosine only & 0.750 & 15/20 & 0.341 \\
Graph (1-hop, $gw$=0.3) & \textbf{1.000} & \textbf{20/20} & 0.341 \\
\bottomrule
\end{tabular}
\caption{Graph expansion achieves 100\% multi-hop recall with zero degradation
to standard retrieval.}
\label{tab:graph}
\end{table}

\subsection{AMBench: Agent Memory Benchmark (WS-D)}

300 incidents, 10 pattern types, 90 simulated days, 3 systems:

\begin{table}[h]
\centering
\begin{tabular}{lrrr}
\toprule
Metric & RAG & IMI ($rw$=0.10) & IMI ($rw$=0.15) \\
\midrule
M1: Retrieval R@5 & 0.279 & 0.275 & 0.273 \\
M2: Cluster purity & 0.736 & 0.736 & 0.736 \\
M3: Action P@1 & 100\% & 100\% & 100\% \\
M4: Temporal avg age & 41.2d & \textbf{16.8d} & 43.0d \\
\bottomrule
\end{tabular}
\caption{AMBench results. IMI with $rw$=0.10 dramatically improves temporal
coherence (avg age 16.8 vs 41.2 days) while maintaining retrieval accuracy.}
\label{tab:ambench}
\end{table}

\subsection{Adaptive Relevance Weight (P1)}

Keyword-based query intent classifier that adjusts $rw$ dynamically:

\begin{table}[h]
\centering
\begin{tabular}{llrl}
\toprule
Intent & Pattern & $rw$ & Example \\
\midrule
Temporal & ``recent'', ``latest'' & 0.15 & ``recent auth failures'' \\
Exploratory & ``find all'', ``list every'' & 0.00 & ``find all cert incidents'' \\
Action & ``how to'', ``fix'' & 0.05 & ``how to prevent DNS'' \\
Default & everything else & 0.10 & ``auth token failures'' \\
\bottomrule
\end{tabular}
\caption{Adaptive relevance weight. Classification accuracy: 100\% (16/16).
MRR: adaptive (0.651) $>$ best fixed (0.643). Zero cost.}
\label{tab:adaptive}
\end{table}

\subsection{Causal Edge Detection (P2)}

Empirical finding: causal pairs have \textbf{low} cosine similarity (avg 0.308).
``Token race condition'' and ``rolling deploy'' are semantically distant---the
relationship is logical, not semantic.

\begin{table}[h]
\centering
\begin{tabular}{lrrl}
\toprule
Strategy & Recall & Precision & Cost \\
\midrule
Embedding ($\theta$=0.40) & 30\% & 10\% & 0 LLM calls \\
Embedding ($\theta$=0.65) & 10\% & 100\% & 0 LLM calls \\
LLM confirmed & $\sim$80\%+ & $\sim$90\%+ & 1 call/pair \\
\bottomrule
\end{tabular}
\caption{Causality detection requires reasoning, not similarity.}
\label{tab:causal}
\end{table}

\section{Architecture Decisions}

Through empirical validation, we made five key decisions:

\begin{enumerate}
  \item \textbf{Removed TimescaleDB backend} ($-$425 lines). SQLite provides
    equivalent functionality at 87$\times$ lower latency.
  \item \textbf{Made predictive coding opt-in}. Surprise boost adds only
    +0.003 to R@5 but costs 2 LLM calls per encode.
  \item \textbf{Set default $rw$ to 0.10} (from 0.30). Ablation showed 0.3
    was too aggressive.
  \item \textbf{Added adaptive $rw$}. Query intent classifier adjusts $rw$
    per query with zero cost, improving MRR by +1.2\%.
  \item \textbf{Causal detection requires LLM}. Embedding similarity is
    insufficient (avg causal pair $\cos = 0.308$).
\end{enumerate}

\section{Limitations}

\begin{itemize}
  \item \textbf{Synthetic evaluation only}: All experiments use generated incidents,
    not real-world agent traces.
  \item \textbf{Single embedding model}: All experiments use all-MiniLM-L6-v2 (384d).
    Larger models may shift the balance.
  \item \textbf{No real HippoRAG comparison}: Compared against regex-NER simulation.
  \item \textbf{Static ground truth}: Relevance judgments are author-created, not
    crowdsourced.
  \item \textbf{Single-agent only}: AMBench simulates one agent.
\end{itemize}

\section{Future Work}

\begin{itemize}
  \item \textbf{Real-world validation}: Deploy IMI in production agent systems.
  \item \textbf{Multi-agent memory sharing}: Agent-scoped views with trust gradients.
  \item \textbf{LLM-confirmed causal detection}: Embedding candidates filtered
    by LLM confirmation.
  \item \textbf{Larger benchmarks}: 10K incidents, 365 days, multi-agent.
\end{itemize}

\section{Conclusion}

IMI demonstrates that cognitive memory features---despite degrading pure
retrieval benchmarks---significantly improve agent-relevant metrics:
temporal coherence (+147\% recency), multi-hop recall (+33\%), and domain
precision (+6.7\%). The optimal configuration uses moderate relevance weighting
($rw$=0.10), graph-augmented expansion for multi-hop, and affective modulation
for emotional salience.

AMBench provides the first standardized evaluation framework for agent memory,
measuring the full encode-retrieve-consolidate-act lifecycle that existing
RAG benchmarks ignore.

IMI is open-source, runs on SQLite with zero infrastructure, requires no LLM
calls at query time, and achieves competitive retrieval quality while adding
temporal, affective, and action-oriented capabilities that pure vector stores
lack.

\section*{Acknowledgments}

This work was developed using Claude Code (Anthropic) as a research
acceleration tool. All experiments, ablations, and benchmark designs were
conceived and validated by the author; Claude Code assisted with
implementation and systematic exploration of the experimental parameter space.

\bibliographystyle{plainnat}

\begin{thebibliography}{7}

\bibitem[Collins and Loftus(1975)]{collins1975}
Collins, A.M. and Loftus, E.F. (1975).
\newblock A spreading-activation theory of semantic processing.
\newblock \emph{Psychological Review}, 82(6):407--428.

\bibitem[Friston(2005)]{friston2005}
Friston, K. (2005).
\newblock A theory of cortical responses.
\newblock \emph{Phil. Trans. R. Soc. B}, 360(1456):815--836.

\bibitem[Gibson(1979)]{gibson1979}
Gibson, J.J. (1979).
\newblock \emph{The Ecological Approach to Visual Perception}.
\newblock Houghton Mifflin.

\bibitem[Gutierrez et~al.(2024)]{gutierrez2024}
Gutierrez, B. et~al. (2024).
\newblock HippoRAG: Neurobiologically Inspired Long-Term Memory for Large
  Language Models.
\newblock In \emph{NeurIPS 2024}.

\bibitem[Packer et~al.(2023)]{packer2023}
Packer, C. et~al. (2023).
\newblock MemGPT: Towards LLMs as Operating Systems.
\newblock \emph{arXiv:2310.08560}.

\bibitem[Rao and Ballard(1999)]{rao1999}
Rao, R.P.N. and Ballard, D.H. (1999).
\newblock Predictive coding in the visual cortex: a functional interpretation
  of some extra-classical receptive-field effects.
\newblock \emph{Nature Neuroscience}, 2(1):79--87.

\bibitem[Sarthi et~al.(2024)]{sarthi2024}
Sarthi, P. et~al. (2024).
\newblock RAPTOR: Recursive Abstractive Processing for Tree-Organized
  Retrieval.
\newblock In \emph{ICLR 2024}.

\end{thebibliography}

\end{document}
